\documentclass{article}

\usepackage[T1]{fontenc}

\usepackage{booktabs,graphicx,siunitx}

\begin{document}

\section{Difference-in-differences (DiD)}

\begin{tabular}{ll}
\toprule
 & (1) \\
\midrule
Post & 2.02 \\
 & (0.09) \\
 & [1.84, 2.19] \\
Treated $\times$ Post & 1.53 \\
 & (0.12) \\
 & [1.29, 1.76] \\
\midrule
Num.Obs. & 96 \\
N entities & 12 \\
Within R\textsuperscript{2} & .84 \\
F & F(2, 82) = 216.31, p $<$ .001 \\
\bottomrule
\end{tabular}

\begin{tabular}{lll}
\toprule
Group & Pre & Post \\
\midrule
Control & 19.21 & 21.22 \\
Treated & 12.92 & 16.46 \\
\bottomrule
\end{tabular}

\includegraphics{figures/01_did/figure_parallel-trends.png}

The Treated$\times$Post coefficient is the estimated treatment effect. It rests on the parallel-trends assumption: that the groups would have moved together absent treatment. Similar pre-treatment trends (inspect the pre-period of the plot) make this more plausible but do not prove it --- a visual check is supportive, not confirmatory, since the assumption is about the unobservable post-period counterfactual. The note adds a formal pre-trends test (a pre-period leads-and-lags joint F of the treated$\times$time interactions): a small p flags diverging pre-trends, while a large p is consistent with parallel trends. Method: R plm package with clustered standard errors (Croissant \& Millo, 2008; Bertrand, Duflo \& Mullainathan, 2004). Your significance threshold ($\alpha$) is 0.05.

The DiD estimate was B=1.53, 95\% CI [1.29, 1.76], p $<$ .001 (clustered SE).

Statistical basis: Difference-in-differences via entity fixed effects (plm within estimator) (Croissant, Y., Millo, G. (2008). "Panel Data Econometrics in R: The plm Package." Journal of Statistical Software, 27(2), 1-43.); Clustered standard errors / parallel-trends interpretation (Bertrand, M., Duflo, E., Mullainathan, S. (2004). "How Much Should We Trust Differences-in-Differences Estimates?" Quarterly Journal of Economics, 119(1), 249-275.).

\end{document}